\documentclass[11pt,a4paper]{article}

\def\courseTitle{Industrial Security (Bachelor)}
\def\labNumber{06}
\def\labTitle{Real CVE Walk-Through -- Solutions}
\def\labSection{06}
\def\labTime{45 min}

% =====================================================================
% Shared preamble for lab sheets.
%
% Caller must \documentclass{article} first, then define:
%   \def\courseTitle{...}
%   \def\labNumber{NN}
%   \def\labTitle{...}
%   \def\labTime{e.g. "30--45 min"}
%   \def\labSection{e.g. "02"}
% Then % =====================================================================
% Shared preamble for lab sheets.
%
% Caller must \documentclass{article} first, then define:
%   \def\courseTitle{...}
%   \def\labNumber{NN}
%   \def\labTitle{...}
%   \def\labTime{e.g. "30--45 min"}
%   \def\labSection{e.g. "02"}
% Then % =====================================================================
% Shared preamble for lab sheets.
%
% Caller must \documentclass{article} first, then define:
%   \def\courseTitle{...}
%   \def\labNumber{NN}
%   \def\labTitle{...}
%   \def\labTime{e.g. "30--45 min"}
%   \def\labSection{e.g. "02"}
% Then \input{../../template/lab-preamble.tex}
% =====================================================================

\usepackage[utf8]{inputenc}
\usepackage[T1]{fontenc}
\usepackage[english]{babel}
\usepackage[a4paper, margin=2cm]{geometry}
\usepackage{graphicx}
\usepackage{listings}
\usepackage{xcolor}
\usepackage{colortbl}
\usepackage{tikz}
\usepackage{amsmath}
\usepackage{amssymb}
\usepackage{booktabs}
\usepackage{enumitem}
\usepackage{titlesec}
\usepackage{fancyhdr}
\usepackage{hyperref}
\usepackage{tcolorbox}
\tcbuselibrary{skins, breakable}
\usepackage{fontawesome5}

\input{../../../template/tikz-styles.tex}

\providecommand{\courseAuthor}{Industrial Security Lecture Series}
\providecommand{\courseInstitute}{Department of Computer Science}
\providecommand{\companyLogo}{logo}

\definecolor{slideAccent}{HTML}{00205B}
\definecolor{slideSecondary}{HTML}{00A0DC}

% --- Corporate branding overrides ------------------------------------
\IfFileExists{../../../branding.tex}{\input{../../../branding.tex}}{}

\colorlet{labAccent}{slideAccent}
\colorlet{labSec}{slideSecondary}
\definecolor{labCheck}{HTML}{2E7D32}
\definecolor{labWarn}{HTML}{B23A48}

\lstset{
  backgroundcolor=\color{black!4},
  basicstyle=\ttfamily\footnotesize,
  breaklines=true,
  frame=leftline,
  framerule=2pt,
  rulecolor=\color{labAccent},
  numbers=left,
  numberstyle=\tiny\color{black!50},
  showstringspaces=false,
  tabsize=2
}

\setlength{\tabcolsep}{4pt}

% Let TeX add a little extra inter-word stretch as a last resort before
% it reports an Overfull \hbox. Only kicks in on lines that would
% otherwise overflow (long \texttt paths, URLs, CIDRs); never loosens a
% line that already fits. Keeps prose/itemize within the margin.
\setlength{\emergencystretch}{3em}

\pagestyle{fancy}
\fancyhf{}
\renewcommand{\headrulewidth}{0.4pt}
\lhead{\small \courseTitle{} -- Lab \labNumber}
\rhead{\small Maps to Section \labSection}
\cfoot{\thepage}

\newtcolorbox{stepbox}[1]{
  enhanced, breakable, arc=2pt, boxrule=0pt,
  colback=labAccent!7, colframe=labAccent, leftrule=3pt,
  fonttitle=\bfseries\small, coltitle=white,
  colbacktitle=labAccent,
  title={\faTools\ Step #1}
}

\newtcolorbox{warnbox}{
  enhanced, breakable, arc=2pt, boxrule=0pt,
  colback=labWarn!7, colframe=labWarn, leftrule=4pt,
  fonttitle=\bfseries\small, coltitle=white, colbacktitle=labWarn,
  title={\faExclamationTriangle\ Caution}
}

\newtcolorbox{notebox}{
  enhanced, breakable, arc=2pt, boxrule=0pt,
  colback=labAccent!4, colframe=labAccent, leftrule=2pt,
  fonttitle=\bfseries\small, coltitle=white, colbacktitle=labAccent,
  title={\faInfoCircle\ Note}
}

\newtcolorbox{goalbox}{
  enhanced, breakable, arc=2pt, boxrule=0pt,
  colback=labCheck!7, colframe=labCheck, leftrule=3pt,
  fonttitle=\bfseries\small, coltitle=white, colbacktitle=labCheck,
  title={\faBullseye\ Goal}
}

\newtcolorbox{deliverbox}{
  enhanced, breakable, arc=2pt, boxrule=0pt,
  colback=labSec!10, colframe=labSec, leftrule=3pt,
  fonttitle=\bfseries\small, coltitle=white, colbacktitle=labSec,
  title={\faClipboardList\ Deliverable}
}

% --- Solution-sheet boxes (used only by lab-solutions.tex) ----------
\newtcolorbox{solutionbox}[1]{
  enhanced, breakable, arc=2pt, boxrule=0pt,
  colback=labCheck!7, colframe=labCheck, leftrule=3pt,
  fonttitle=\bfseries\small, coltitle=white, colbacktitle=labCheck,
  title={\faCheckCircle\ #1}
}

\newtcolorbox{rubricbox}{
  enhanced, breakable, arc=2pt, boxrule=0pt,
  colback=labSec!7, colframe=labSec, leftrule=3pt,
  fonttitle=\bfseries\small, coltitle=white, colbacktitle=labSec,
  title={\faBalanceScale\ Grading rubric}
}

\titleformat{\section}{\large\bfseries\color{labAccent}}{\thesection}{0.5em}{}

\newcommand{\labheader}{%
  \noindent
  \begin{tikzpicture}
    \fill[labAccent] (0,0) rectangle (\textwidth, -1.6cm);
    \node[anchor=north west, text=white, font=\bfseries\large] at (0.6cm, -0.4cm) {Lab \labNumber{} -- \labTitle};
    \node[anchor=south west, text=white, font=\small] at (0.6cm, -1.3cm) {\courseTitle{} \quad $\bullet$ \quad Maps to Section \labSection{} \quad $\bullet$ \quad \labTime};
    \node[anchor=north east, fill=labSec, text=white, font=\bfseries, inner sep=4pt, rounded corners=2pt]
      at ([xshift=-0.4cm, yshift=-0.4cm]\textwidth,0) {LAB \labNumber};
  \end{tikzpicture}
  \vspace{4mm}
}

% =====================================================================

\usepackage[utf8]{inputenc}
\usepackage[T1]{fontenc}
\usepackage[english]{babel}
\usepackage[a4paper, margin=2cm]{geometry}
\usepackage{graphicx}
\usepackage{listings}
\usepackage{xcolor}
\usepackage{colortbl}
\usepackage{tikz}
\usepackage{amsmath}
\usepackage{amssymb}
\usepackage{booktabs}
\usepackage{enumitem}
\usepackage{titlesec}
\usepackage{fancyhdr}
\usepackage{hyperref}
\usepackage{tcolorbox}
\tcbuselibrary{skins, breakable}
\usepackage{fontawesome5}

% =====================================================================
% Shared TikZ styles for the Industrial Security lecture series.
% Loaded by every slide deck and exercise sheet that uses TikZ.
% =====================================================================

\usetikzlibrary{
  arrows.meta,
  shapes,
  shapes.geometric,
  shapes.symbols,
  shapes.multipart,
  positioning,
  fit,
  calc,
  backgrounds,
  decorations.pathmorphing,
  decorations.pathreplacing,
  decorations.text,
  patterns,
  matrix,
  shadows
}

% --- colour palette --------------------------------------------------
\definecolor{isOT}{HTML}{1F4E79}     % OT (deep blue)
\definecolor{isIT}{HTML}{6F2DA8}     % IT (purple)
\definecolor{isSafety}{HTML}{C00000} % safety red
\definecolor{isNet}{HTML}{2E7D32}    % network green
\definecolor{isAttack}{HTML}{B23A48} % attack red
\definecolor{isAccent}{HTML}{E08E0B} % amber
\definecolor{isMute}{HTML}{6B6B6B}   % grey
\definecolor{isLight}{HTML}{EDF2F8}  % light background
\definecolor{isPLC}{HTML}{2C5282}
\definecolor{isHMI}{HTML}{2F855A}
\definecolor{isSCADA}{HTML}{6B46C1}

% --- generic node styles --------------------------------------------
\tikzset{
  isnode/.style={
    draw, rounded corners=2pt, minimum height=8mm, minimum width=18mm,
    font=\small, align=center, thick
  },
  isbox/.style={isnode, fill=isLight},
  plc/.style={isnode, fill=isPLC!15, draw=isPLC, thick},
  hmi/.style={isnode, fill=isHMI!15, draw=isHMI, thick},
  scada/.style={isnode, fill=isSCADA!15, draw=isSCADA, thick},
  sensor/.style={isnode, fill=isNet!10, draw=isNet, minimum height=6mm, minimum width=14mm, font=\scriptsize},
  actuator/.style={isnode, fill=isAccent!15, draw=isAccent, minimum height=6mm, minimum width=14mm, font=\scriptsize},
  firewall/.style={isnode, fill=isAttack!15, draw=isAttack, thick},
  server/.style={isnode, fill=isMute!15, draw=isMute!80, thick},
  switch/.style={isnode, fill=isLight, draw=black!60, minimum height=6mm, minimum width=14mm, font=\scriptsize},
  workstation/.style={isnode, fill=white, draw=isOT, thick},
  attacker/.style={isnode, fill=isAttack!20, draw=isAttack, thick, font=\small\bfseries},
  inetnode/.style={draw, ellipse, minimum width=24mm, minimum height=10mm, font=\scriptsize, fill=isLight, thick},
  process/.style={isnode, fill=isOT!15, draw=isOT, thick, minimum width=24mm},
  zone/.style={
    draw, rounded corners=4pt, thick, dashed, inner sep=4mm
  },
  conduit/.style={-{Stealth[length=2.5mm]}, thick, isNet!80!black},
  attack/.style={-{Stealth[length=2.5mm]}, thick, isAttack, dashed},
  flow/.style={-{Stealth[length=2.5mm]}, thick, isOT!70!black},
  netline/.style={thick, black!60},
  axisline/.style={-{Stealth[length=2mm]}, thick, black!70},
  tag/.style={font=\scriptsize\itshape, fill=white, inner sep=1pt},
  level/.style={
    draw, fill=isLight, minimum width=0.85\linewidth, minimum height=8mm,
    font=\small, anchor=west
  },
  brace/.style={decorate, decoration={brace, amplitude=4pt, raise=2pt}},
  legendbox/.style={draw, rounded corners=2pt, fill=isLight, inner sep=2mm, font=\scriptsize, align=left}
}

% --- handy macros ----------------------------------------------------
% \plcicon, \hmiicon etc as simple labelled boxes; useful inline.
\newcommand{\plcicon}[1][PLC]{\tikz[baseline=-0.6ex]{\node[plc, minimum width=10mm, minimum height=5mm, font=\scriptsize, inner sep=1pt]{#1};}}
\newcommand{\hmiicon}[1][HMI]{\tikz[baseline=-0.6ex]{\node[hmi, minimum width=10mm, minimum height=5mm, font=\scriptsize, inner sep=1pt]{#1};}}
\newcommand{\fwicon}[1][FW]{\tikz[baseline=-0.6ex]{\node[firewall, minimum width=10mm, minimum height=5mm, font=\scriptsize, inner sep=1pt]{#1};}}


\providecommand{\courseAuthor}{Industrial Security Lecture Series}
\providecommand{\courseInstitute}{Department of Computer Science}
\providecommand{\companyLogo}{logo}

\definecolor{slideAccent}{HTML}{00205B}
\definecolor{slideSecondary}{HTML}{00A0DC}

% --- Corporate branding overrides ------------------------------------
\IfFileExists{../../../branding.tex}{% =====================================================================
% Corporate / institutional branding -- single file to customise the
% look of the entire course (slides, notes, exercises, labs).
%
% HOW TO REBRAND IN UNDER A MINUTE:
%   1. Replace `branding/logo.pdf` with your own logo
%      (PDF preferred; PNG and JPG also work -- name it `logo.<ext>`).
%   2. (Optional) change the two colours below to your brand palette.
%   3. (Optional) override the author/institute lines below.
%   4. Re-run `make` at the repo root. That's it.
%
% Everything in this file has sensible defaults; you can delete a line
% to fall back to the built-in defaults, or comment it out.
%
% This file is loaded automatically by the slides, exercise, and lab
% preambles -- you never need to \input it manually.
% =====================================================================

% --- Logo --------------------------------------------------------------
% Path is resolved from each leaf-build directory; the default points at
% the shared `branding/` folder at the repo root. If you put your logo
% somewhere else, set the full or relative path here (without extension):
\renewcommand{\companyLogo}{../../../branding/logo}

% --- Author / institute (appears on the title page footer) ------------
\renewcommand{\courseAuthor}{}
\renewcommand{\courseInstitute}{}

% --- Brand colours ----------------------------------------------------
% slideAccent      = primary brand colour (title bands, accent rule)
% slideSecondary   = secondary brand colour (section badge, highlight)
% Both use HTML hex codes. Keep enough contrast against white text.
\definecolor{slideAccent}{HTML}{00205B}      % deep navy
\definecolor{slideSecondary}{HTML}{00A0DC}   % light blue accent
}{}

\colorlet{labAccent}{slideAccent}
\colorlet{labSec}{slideSecondary}
\definecolor{labCheck}{HTML}{2E7D32}
\definecolor{labWarn}{HTML}{B23A48}

\lstset{
  backgroundcolor=\color{black!4},
  basicstyle=\ttfamily\footnotesize,
  breaklines=true,
  frame=leftline,
  framerule=2pt,
  rulecolor=\color{labAccent},
  numbers=left,
  numberstyle=\tiny\color{black!50},
  showstringspaces=false,
  tabsize=2
}

\setlength{\tabcolsep}{4pt}

% Let TeX add a little extra inter-word stretch as a last resort before
% it reports an Overfull \hbox. Only kicks in on lines that would
% otherwise overflow (long \texttt paths, URLs, CIDRs); never loosens a
% line that already fits. Keeps prose/itemize within the margin.
\setlength{\emergencystretch}{3em}

\pagestyle{fancy}
\fancyhf{}
\renewcommand{\headrulewidth}{0.4pt}
\lhead{\small \courseTitle{} -- Lab \labNumber}
\rhead{\small Maps to Section \labSection}
\cfoot{\thepage}

\newtcolorbox{stepbox}[1]{
  enhanced, breakable, arc=2pt, boxrule=0pt,
  colback=labAccent!7, colframe=labAccent, leftrule=3pt,
  fonttitle=\bfseries\small, coltitle=white,
  colbacktitle=labAccent,
  title={\faTools\ Step #1}
}

\newtcolorbox{warnbox}{
  enhanced, breakable, arc=2pt, boxrule=0pt,
  colback=labWarn!7, colframe=labWarn, leftrule=4pt,
  fonttitle=\bfseries\small, coltitle=white, colbacktitle=labWarn,
  title={\faExclamationTriangle\ Caution}
}

\newtcolorbox{notebox}{
  enhanced, breakable, arc=2pt, boxrule=0pt,
  colback=labAccent!4, colframe=labAccent, leftrule=2pt,
  fonttitle=\bfseries\small, coltitle=white, colbacktitle=labAccent,
  title={\faInfoCircle\ Note}
}

\newtcolorbox{goalbox}{
  enhanced, breakable, arc=2pt, boxrule=0pt,
  colback=labCheck!7, colframe=labCheck, leftrule=3pt,
  fonttitle=\bfseries\small, coltitle=white, colbacktitle=labCheck,
  title={\faBullseye\ Goal}
}

\newtcolorbox{deliverbox}{
  enhanced, breakable, arc=2pt, boxrule=0pt,
  colback=labSec!10, colframe=labSec, leftrule=3pt,
  fonttitle=\bfseries\small, coltitle=white, colbacktitle=labSec,
  title={\faClipboardList\ Deliverable}
}

% --- Solution-sheet boxes (used only by lab-solutions.tex) ----------
\newtcolorbox{solutionbox}[1]{
  enhanced, breakable, arc=2pt, boxrule=0pt,
  colback=labCheck!7, colframe=labCheck, leftrule=3pt,
  fonttitle=\bfseries\small, coltitle=white, colbacktitle=labCheck,
  title={\faCheckCircle\ #1}
}

\newtcolorbox{rubricbox}{
  enhanced, breakable, arc=2pt, boxrule=0pt,
  colback=labSec!7, colframe=labSec, leftrule=3pt,
  fonttitle=\bfseries\small, coltitle=white, colbacktitle=labSec,
  title={\faBalanceScale\ Grading rubric}
}

\titleformat{\section}{\large\bfseries\color{labAccent}}{\thesection}{0.5em}{}

\newcommand{\labheader}{%
  \noindent
  \begin{tikzpicture}
    \fill[labAccent] (0,0) rectangle (\textwidth, -1.6cm);
    \node[anchor=north west, text=white, font=\bfseries\large] at (0.6cm, -0.4cm) {Lab \labNumber{} -- \labTitle};
    \node[anchor=south west, text=white, font=\small] at (0.6cm, -1.3cm) {\courseTitle{} \quad $\bullet$ \quad Maps to Section \labSection{} \quad $\bullet$ \quad \labTime};
    \node[anchor=north east, fill=labSec, text=white, font=\bfseries, inner sep=4pt, rounded corners=2pt]
      at ([xshift=-0.4cm, yshift=-0.4cm]\textwidth,0) {LAB \labNumber};
  \end{tikzpicture}
  \vspace{4mm}
}

% =====================================================================

\usepackage[utf8]{inputenc}
\usepackage[T1]{fontenc}
\usepackage[english]{babel}
\usepackage[a4paper, margin=2cm]{geometry}
\usepackage{graphicx}
\usepackage{listings}
\usepackage{xcolor}
\usepackage{colortbl}
\usepackage{tikz}
\usepackage{amsmath}
\usepackage{amssymb}
\usepackage{booktabs}
\usepackage{enumitem}
\usepackage{titlesec}
\usepackage{fancyhdr}
\usepackage{hyperref}
\usepackage{tcolorbox}
\tcbuselibrary{skins, breakable}
\usepackage{fontawesome5}

% =====================================================================
% Shared TikZ styles for the Industrial Security lecture series.
% Loaded by every slide deck and exercise sheet that uses TikZ.
% =====================================================================

\usetikzlibrary{
  arrows.meta,
  shapes,
  shapes.geometric,
  shapes.symbols,
  shapes.multipart,
  positioning,
  fit,
  calc,
  backgrounds,
  decorations.pathmorphing,
  decorations.pathreplacing,
  decorations.text,
  patterns,
  matrix,
  shadows
}

% --- colour palette --------------------------------------------------
\definecolor{isOT}{HTML}{1F4E79}     % OT (deep blue)
\definecolor{isIT}{HTML}{6F2DA8}     % IT (purple)
\definecolor{isSafety}{HTML}{C00000} % safety red
\definecolor{isNet}{HTML}{2E7D32}    % network green
\definecolor{isAttack}{HTML}{B23A48} % attack red
\definecolor{isAccent}{HTML}{E08E0B} % amber
\definecolor{isMute}{HTML}{6B6B6B}   % grey
\definecolor{isLight}{HTML}{EDF2F8}  % light background
\definecolor{isPLC}{HTML}{2C5282}
\definecolor{isHMI}{HTML}{2F855A}
\definecolor{isSCADA}{HTML}{6B46C1}

% --- generic node styles --------------------------------------------
\tikzset{
  isnode/.style={
    draw, rounded corners=2pt, minimum height=8mm, minimum width=18mm,
    font=\small, align=center, thick
  },
  isbox/.style={isnode, fill=isLight},
  plc/.style={isnode, fill=isPLC!15, draw=isPLC, thick},
  hmi/.style={isnode, fill=isHMI!15, draw=isHMI, thick},
  scada/.style={isnode, fill=isSCADA!15, draw=isSCADA, thick},
  sensor/.style={isnode, fill=isNet!10, draw=isNet, minimum height=6mm, minimum width=14mm, font=\scriptsize},
  actuator/.style={isnode, fill=isAccent!15, draw=isAccent, minimum height=6mm, minimum width=14mm, font=\scriptsize},
  firewall/.style={isnode, fill=isAttack!15, draw=isAttack, thick},
  server/.style={isnode, fill=isMute!15, draw=isMute!80, thick},
  switch/.style={isnode, fill=isLight, draw=black!60, minimum height=6mm, minimum width=14mm, font=\scriptsize},
  workstation/.style={isnode, fill=white, draw=isOT, thick},
  attacker/.style={isnode, fill=isAttack!20, draw=isAttack, thick, font=\small\bfseries},
  inetnode/.style={draw, ellipse, minimum width=24mm, minimum height=10mm, font=\scriptsize, fill=isLight, thick},
  process/.style={isnode, fill=isOT!15, draw=isOT, thick, minimum width=24mm},
  zone/.style={
    draw, rounded corners=4pt, thick, dashed, inner sep=4mm
  },
  conduit/.style={-{Stealth[length=2.5mm]}, thick, isNet!80!black},
  attack/.style={-{Stealth[length=2.5mm]}, thick, isAttack, dashed},
  flow/.style={-{Stealth[length=2.5mm]}, thick, isOT!70!black},
  netline/.style={thick, black!60},
  axisline/.style={-{Stealth[length=2mm]}, thick, black!70},
  tag/.style={font=\scriptsize\itshape, fill=white, inner sep=1pt},
  level/.style={
    draw, fill=isLight, minimum width=0.85\linewidth, minimum height=8mm,
    font=\small, anchor=west
  },
  brace/.style={decorate, decoration={brace, amplitude=4pt, raise=2pt}},
  legendbox/.style={draw, rounded corners=2pt, fill=isLight, inner sep=2mm, font=\scriptsize, align=left}
}

% --- handy macros ----------------------------------------------------
% \plcicon, \hmiicon etc as simple labelled boxes; useful inline.
\newcommand{\plcicon}[1][PLC]{\tikz[baseline=-0.6ex]{\node[plc, minimum width=10mm, minimum height=5mm, font=\scriptsize, inner sep=1pt]{#1};}}
\newcommand{\hmiicon}[1][HMI]{\tikz[baseline=-0.6ex]{\node[hmi, minimum width=10mm, minimum height=5mm, font=\scriptsize, inner sep=1pt]{#1};}}
\newcommand{\fwicon}[1][FW]{\tikz[baseline=-0.6ex]{\node[firewall, minimum width=10mm, minimum height=5mm, font=\scriptsize, inner sep=1pt]{#1};}}


\providecommand{\courseAuthor}{Industrial Security Lecture Series}
\providecommand{\courseInstitute}{Department of Computer Science}
\providecommand{\companyLogo}{logo}

\definecolor{slideAccent}{HTML}{00205B}
\definecolor{slideSecondary}{HTML}{00A0DC}

% --- Corporate branding overrides ------------------------------------
\IfFileExists{../../../branding.tex}{% =====================================================================
% Corporate / institutional branding -- single file to customise the
% look of the entire course (slides, notes, exercises, labs).
%
% HOW TO REBRAND IN UNDER A MINUTE:
%   1. Replace `branding/logo.pdf` with your own logo
%      (PDF preferred; PNG and JPG also work -- name it `logo.<ext>`).
%   2. (Optional) change the two colours below to your brand palette.
%   3. (Optional) override the author/institute lines below.
%   4. Re-run `make` at the repo root. That's it.
%
% Everything in this file has sensible defaults; you can delete a line
% to fall back to the built-in defaults, or comment it out.
%
% This file is loaded automatically by the slides, exercise, and lab
% preambles -- you never need to \input it manually.
% =====================================================================

% --- Logo --------------------------------------------------------------
% Path is resolved from each leaf-build directory; the default points at
% the shared `branding/` folder at the repo root. If you put your logo
% somewhere else, set the full or relative path here (without extension):
\renewcommand{\companyLogo}{../../../branding/logo}

% --- Author / institute (appears on the title page footer) ------------
\renewcommand{\courseAuthor}{}
\renewcommand{\courseInstitute}{}

% --- Brand colours ----------------------------------------------------
% slideAccent      = primary brand colour (title bands, accent rule)
% slideSecondary   = secondary brand colour (section badge, highlight)
% Both use HTML hex codes. Keep enough contrast against white text.
\definecolor{slideAccent}{HTML}{00205B}      % deep navy
\definecolor{slideSecondary}{HTML}{00A0DC}   % light blue accent
}{}

\colorlet{labAccent}{slideAccent}
\colorlet{labSec}{slideSecondary}
\definecolor{labCheck}{HTML}{2E7D32}
\definecolor{labWarn}{HTML}{B23A48}

\lstset{
  backgroundcolor=\color{black!4},
  basicstyle=\ttfamily\footnotesize,
  breaklines=true,
  frame=leftline,
  framerule=2pt,
  rulecolor=\color{labAccent},
  numbers=left,
  numberstyle=\tiny\color{black!50},
  showstringspaces=false,
  tabsize=2
}

\setlength{\tabcolsep}{4pt}

% Let TeX add a little extra inter-word stretch as a last resort before
% it reports an Overfull \hbox. Only kicks in on lines that would
% otherwise overflow (long \texttt paths, URLs, CIDRs); never loosens a
% line that already fits. Keeps prose/itemize within the margin.
\setlength{\emergencystretch}{3em}

\pagestyle{fancy}
\fancyhf{}
\renewcommand{\headrulewidth}{0.4pt}
\lhead{\small \courseTitle{} -- Lab \labNumber}
\rhead{\small Maps to Section \labSection}
\cfoot{\thepage}

\newtcolorbox{stepbox}[1]{
  enhanced, breakable, arc=2pt, boxrule=0pt,
  colback=labAccent!7, colframe=labAccent, leftrule=3pt,
  fonttitle=\bfseries\small, coltitle=white,
  colbacktitle=labAccent,
  title={\faTools\ Step #1}
}

\newtcolorbox{warnbox}{
  enhanced, breakable, arc=2pt, boxrule=0pt,
  colback=labWarn!7, colframe=labWarn, leftrule=4pt,
  fonttitle=\bfseries\small, coltitle=white, colbacktitle=labWarn,
  title={\faExclamationTriangle\ Caution}
}

\newtcolorbox{notebox}{
  enhanced, breakable, arc=2pt, boxrule=0pt,
  colback=labAccent!4, colframe=labAccent, leftrule=2pt,
  fonttitle=\bfseries\small, coltitle=white, colbacktitle=labAccent,
  title={\faInfoCircle\ Note}
}

\newtcolorbox{goalbox}{
  enhanced, breakable, arc=2pt, boxrule=0pt,
  colback=labCheck!7, colframe=labCheck, leftrule=3pt,
  fonttitle=\bfseries\small, coltitle=white, colbacktitle=labCheck,
  title={\faBullseye\ Goal}
}

\newtcolorbox{deliverbox}{
  enhanced, breakable, arc=2pt, boxrule=0pt,
  colback=labSec!10, colframe=labSec, leftrule=3pt,
  fonttitle=\bfseries\small, coltitle=white, colbacktitle=labSec,
  title={\faClipboardList\ Deliverable}
}

% --- Solution-sheet boxes (used only by lab-solutions.tex) ----------
\newtcolorbox{solutionbox}[1]{
  enhanced, breakable, arc=2pt, boxrule=0pt,
  colback=labCheck!7, colframe=labCheck, leftrule=3pt,
  fonttitle=\bfseries\small, coltitle=white, colbacktitle=labCheck,
  title={\faCheckCircle\ #1}
}

\newtcolorbox{rubricbox}{
  enhanced, breakable, arc=2pt, boxrule=0pt,
  colback=labSec!7, colframe=labSec, leftrule=3pt,
  fonttitle=\bfseries\small, coltitle=white, colbacktitle=labSec,
  title={\faBalanceScale\ Grading rubric}
}

\titleformat{\section}{\large\bfseries\color{labAccent}}{\thesection}{0.5em}{}

\newcommand{\labheader}{%
  \noindent
  \begin{tikzpicture}
    \fill[labAccent] (0,0) rectangle (\textwidth, -1.6cm);
    \node[anchor=north west, text=white, font=\bfseries\large] at (0.6cm, -0.4cm) {Lab \labNumber{} -- \labTitle};
    \node[anchor=south west, text=white, font=\small] at (0.6cm, -1.3cm) {\courseTitle{} \quad $\bullet$ \quad Maps to Section \labSection{} \quad $\bullet$ \quad \labTime};
    \node[anchor=north east, fill=labSec, text=white, font=\bfseries, inner sep=4pt, rounded corners=2pt]
      at ([xshift=-0.4cm, yshift=-0.4cm]\textwidth,0) {LAB \labNumber};
  \end{tikzpicture}
  \vspace{4mm}
}


\begin{document}
\labheader

\begin{goalbox}
Lecturer / TA reference for Lab 06: a fully worked CISA advisory decomposition, an environmental-CVSS re-scoring example, the annotated safe-nmap profile against the lab PLC, the survival check, and a grading rubric.
\end{goalbox}

\begin{solutionbox}{Part 1, Steps 1--2 -- worked advisory: ICSA-22-102-02}
Reference advisory: \textbf{ICSA-22-102-02}, \emph{Rockwell Automation Logix Controllers}, originally published \textbf{12 April 2022} by CISA. Use this as the canonical worked example; any current advisory of comparable depth is acceptable from a student.

\begin{center}\footnotesize
\begin{tabular}{@{}p{3.6cm}p{11.0cm}@{}}
\toprule
\textbf{Field} & \textbf{Value} \\
\midrule
ICSA ID                 & \texttt{ICSA-22-102-02} \\
Publish date            & 12 Apr 2022 (updated thereafter by CISA) \\
Vendor                  & Rockwell Automation \\
Affected products       & Logix Designer (Studio 5000) + ControlLogix 5570/5580, CompactLogix 5370/5380/5480, GuardLogix 5570/5580, Compact GuardLogix 5370/5380 \\
Affected firmware       & Controllers running affected versions prior to the fixed firmware listed by the vendor PSIRT (see Rockwell PN1645) \\
CVE ID                  & \texttt{CVE-2022-1161} (controller side; companion \texttt{CVE-2022-1159} covers the Studio~5000 side) \\
CVSS v3.1 base score    & \texttt{10.0 (Critical)} \\
CVSS v3.1 vector        & \texttt{CVSS:3.1/AV:N/AC:L/PR:N/UI:N/S:C/C:H/I:H/A:H} \\
CWE                     & \texttt{CWE-20} Improper Input Validation (``forged logic download'' -- user-readable logic differs from logic actually executed) \\
Pre-conditions          & Network reachability of the controller's EtherNet/IP service (TCP/UDP \texttt{44818}, TCP \texttt{2222}); no authentication required if the controller is not in \emph{Run-mode key} hardware lock; no user interaction \\
Stated impact           & Attacker can write modified ladder/structured-text logic to the controller while the engineering workstation continues to display the original program -- a stealthy logic-tamper primitive \\
Known exploitation      & Proof-of-concept disclosed by Claroty (Team82); no confirmed in-the-wild abuse at publish date \\
Vendor mitigation       & Apply fixed firmware per Rockwell PSIRT advisory \texttt{PN1645}; set the controller mode switch to \textbf{RUN} (hardware key) to block remote program downloads \\
CISA mitigation         & Minimise network exposure (controllers not Internet-facing), place behind firewalls, isolate from business network, use secure remote access (VPN) and update the VPN, perform impact analysis before deploying \\
\bottomrule
\end{tabular}
\end{center}

{\footnotesize\emph{Sources:} CISA, \emph{ICSA-22-102-02 Rockwell Automation Logix Controllers}, 12 Apr 2022 \textbar{} Rockwell PSIRT \texttt{PN1645}, 2022 \textbar{} Claroty Team82, \emph{Evil PLC \& Forged Logic Download}, 2022.}
\end{solutionbox}

\begin{solutionbox}{Part 1, Step 3 -- environmental re-scoring (PLC behind a data diode)}
Same controller, deployed in a small water utility behind a hardware one-way data diode (process LAN $\rightarrow$ historian only; no inbound path):

\begin{itemize}[leftmargin=*, nosep]
  \item \texttt{MAV:A} -- Modified Attack Vector \emph{Adjacent}. The diode breaks routed L3 reachability; only an attacker already on the local OT segment can talk EtherNet/IP to the controller.
  \item \texttt{MAC:H} -- Modified Attack Complexity \emph{High}. The attacker must first defeat physical/network access controls (cabinet lock, managed switch, 802.1X) to land on that segment.
  \item \texttt{MPR:N} -- no authentication on the protocol itself, unchanged.
  \item \texttt{MUI:N} -- no user interaction needed, unchanged.
  \item \texttt{MS:U} -- Modified Scope \emph{Unchanged}. The diode prevents the compromised controller from being used as a pivot back into IT; the blast radius stops at the OT cell.
  \item \texttt{MC:H / MI:H / MA:H} -- once on-segment, impact on the process is still total.
  \item Environmental requirements (utility is safety-relevant): \texttt{CR:H, IR:H, AR:H}.
\end{itemize}

Modified vector:
\begin{quote}
\ttfamily\small CVSS:3.1/AV:N/AC:L/PR:N/UI:N/S:C/C:H/I:H/A:H/\\MAV:A/MAC:H/MS:U/CR:H/IR:H/AR:H
\end{quote}

Environmental score drops from \textbf{10.0 Critical} to roughly \textbf{7.0--7.5 High} (exact value depends on the calculator; accept anything in the \textbf{6.8--7.6} band with the justifications above). Teaching point: the vulnerability is just as dangerous \emph{if} the attacker gets on-segment -- the diode bought \emph{exposure} reduction, not \emph{impact} reduction.
\end{solutionbox}

\begin{solutionbox}{Part 2 -- safe nmap against the lab PLC}
Exact command students run inside the student container:

\begin{lstlisting}[language=bash]
$ docker compose exec student \
    nmap -sT -Pn -n --max-rate 5 --scan-delay 200ms \
         -p 502,8080 172.28.0.10
\end{lstlisting}

Flag-by-flag justification:
\begin{itemize}[leftmargin=*, nosep]
  \item \texttt{-sT} -- full TCP connect via the OS stack, \emph{not} raw SYN. Fragile PLC TCP stacks (legacy Schneider, older Siemens S7) have crashed on half-open SYN floods; a clean 3-way handshake plus RST is the safest probe.
  \item \texttt{-Pn} -- skip host discovery. Many PLCs and ICS gateways drop or rate-limit ICMP; without \texttt{-Pn} nmap would mark the target ``down'' and never scan it.
  \item \texttt{-n} -- no DNS resolution. Avoids leaking the target IP to a DNS resolver and removes a source of timing jitter.
  \item \texttt{--max-rate 5} -- cap at 5 packets per second. Well below any documented denial-of-service threshold for ICS devices.
  \item \texttt{--scan-delay 200ms} -- force at least 200\,ms between probes. Belt-and-braces with \texttt{--max-rate}; the two together guarantee a gentle scan even if nmap recalculates timing.
  \item \texttt{-p 502,8080} -- only the two ports the lab actually uses (Modbus/TCP and the OpenPLC web UI). No port sweep, no UDP, no fingerprinting probes.
\end{itemize}

Expected output (abridged):
\begin{lstlisting}[language=bash]
Starting Nmap 7.94 ( https://nmap.org ) at 2026-... CET
Nmap scan report for 172.28.0.10
Host is up (0.00050s latency).

PORT     STATE SERVICE
502/tcp  open  mbap
8080/tcp open  http-proxy

Nmap done: 1 IP address (1 host up, 2 ports scanned) in ~1.6s
\end{lstlisting}
\end{solutionbox}

\begin{solutionbox}{Part 2 -- forbidden flags and why}
Do \textbf{not} run any of the following against a real PLC (or this lab, which mimics one):
\begin{itemize}[leftmargin=*, nosep]
  \item \texttt{-sV} (version detection) -- sends protocol-specific probes (\texttt{NULL}, Modbus malformed PDUs, HTTP \texttt{OPTIONS}, etc.). Documented to lock ladder-logic memory regions on several controllers and to crash older Schneider Modicon Quantum CPUs during banner grabs.
  \item \texttt{-O} (OS fingerprint) -- sends crafted TCP/IP packets (out-of-order flags, weird window sizes). Legacy ICS TCP stacks have rebooted under \texttt{-O}.
  \item \texttt{--script} / \texttt{-sC} -- NSE scripts vary wildly in safety; \texttt{modbus-discover}, \texttt{s7-info}, etc.\ have all triggered comms faults on production PLCs.
  \item \texttt{-A} -- shorthand for \texttt{-O -sV -sC --traceroute}: combines everything above. \textbf{Never} use it on OT.
\end{itemize}

{\footnotesize\emph{Sources:} Digital Bond / Project Basecamp, \emph{ICS Device Failure Modes During Active Scanning}, 2012 \textbar{} ICS-CERT, \emph{Recommended Practice: Updating Antivirus in an Industrial Control System}, 2011 -- both document the historical fragility.}
\end{solutionbox}

\begin{solutionbox}{Part 2 -- post-scan survival check}
After the nmap run, the student must re-run the Lab 02 read script:
\begin{lstlisting}[language=bash]
$ docker compose exec student python3 /labs/scripts/modbus_read.py
Registers 0..9: [0, 0, 0, 0, 0, 0, 0, 0, 0, 0]
\end{lstlisting}
A successful read of \textbf{exactly 10 registers} proves the Modbus/TCP server is still responsive and the IEC task is still cycling. If the script throws \texttt{ConnectionRefusedError} or hangs past 5\,s, the PLC stack did not survive the scan -- collect the nmap command actually issued, the container logs (\texttt{docker compose logs openplc}), and stop further scanning until the runtime is restarted.
\end{solutionbox}

\begin{solutionbox}{Common student mistakes}
\begin{itemize}[leftmargin=*, nosep]
  \item \textbf{\texttt{nmap -sS}} -- raw SYN scan; needs \texttt{root}/\texttt{CAP\_NET\_RAW} in the container and is exactly the half-open behaviour that has crashed legacy PLC stacks. Mark down even though the lab container ``survives'' it.
  \item \textbf{Default scan with no rate cap} -- \texttt{nmap 172.28.0.10} alone fires $\sim$1000 probes at line rate; cite that pps figure when grading.
  \item \textbf{Wrong target IP} -- scanning \texttt{127.0.0.1} from inside \texttt{is-student} hits the student container, not OpenPLC. Correct target is \texttt{172.28.0.10}.
  \item \textbf{\texttt{-A} aggressive scan} -- collects everything (\texttt{-O -sV -sC --traceroute}). Auto-fail this part.
  \item \textbf{Skipping the survival check} -- nmap completing successfully is not evidence the PLC is healthy; only the follow-up \texttt{modbus\_read.py} is.
  \item \textbf{Citing the wrong vendor advisory} -- students sometimes paraphrase a news article instead of the CISA ICSA page; require the \texttt{ICSA-YY-NNN-NN} identifier and the vendor PSIRT number.
\end{itemize}
\end{solutionbox}

\begin{rubricbox}
\begin{itemize}[leftmargin=*, nosep]
  \item \textbf{3 pts} -- advisory decomposition complete: ICSA ID, vendor, affected models + firmware band, CVE ID(s), CVSS v3.1 vector \emph{and} score, CWE, vendor mitigation, CISA mitigation.
  \item \textbf{2 pts} -- environmental re-score: at least \texttt{MAV}, \texttt{MAC}, \texttt{MS} modifiers justified in writing; numerical result within $\pm$0.5 of the worked answer.
  \item \textbf{2 pts} -- safe-nmap command run exactly as specified (\texttt{-sT -Pn -n --max-rate 5 --scan-delay 200ms -p 502,8080}) against \texttt{172.28.0.10}; output captured.
  \item \textbf{1 pt} -- post-scan \texttt{modbus\_read.py} returns the expected 10-register vector.
  \item \textbf{2 pts} -- short written paragraph (3--5 sentences) on \emph{why} \texttt{-sV}, \texttt{-O}, \texttt{-A}, \texttt{--script} are unsafe against real PLCs, with at least one concrete vendor/incident reference.
\end{itemize}
Total: \textbf{10 pts}. Auto-deduct 2 pts for any forbidden flag in the submitted command history; auto-fail (0 pts on Part 2) if the student scanned an IP outside the lab subnet \texttt{172.28.0.0/24}.
\end{rubricbox}

\end{document}
